\documentclass[11pt]{article}
\usepackage[margin=1in]{geometry}
\usepackage{amsmath,amssymb,booktabs}
\usepackage{hyperref}

\title{Two-Batch Uplift Monitoring:\\
FSDS-Isomorphic Formulation with LOCO--AUUC and Two-Layer Localization}
\author{Uplift FSDS monitoring note}
\date{}

\begin{document}
\maketitle

\begin{abstract}
We formalise uplift model monitoring under a \textbf{two-batch} contract:
REF calibrates a frozen ranker $\hat\tau$; LIVE is the deployment window.
Batch indicator $W$ aligns with standard FSDS on covariates; treatment $T$ enters AUUC and slice-wise empirical ATE.
Layer~1 (AUUC, LOCO--AUUC, domain-quintile pairwise comparisons) localises \emph{ranking} failure;
Layer~2 (empirical ATE and mean $\hat\tau$ on LIVE slices) localises \emph{effect movement}, gated by PO-risk (DRPerm) on stacked batches.
The pipeline outputs direct targeting policies without a hierarchical feature tree.
\end{abstract}

\noindent\textbf{Overview.}\ 
The post-hoc localization is performed at two levels: one based on AUUC (evaluated on a single batch) and the other based on estimated uplift, followed by pseudo-outcome risk.
This design is highly efficient and beneficial for downstream procedures.
Moreover, the AUC of the RF-domain binary classifier is instrumental for monitoring covariate shift.
These elements collectively produce compound insights that facilitate interpretation.

\medskip
% Canonical LaTeX formulation: two-batch (REF vs LIVE) uplift monitoring
% isomorphic to batch FSDS, with AUUC ranking layer + uplift effect layer.
% Compile via uplift_fsds_two_batch_standalone.tex

\section{Two-batch uplift monitoring (FSDS-isomorphic formulation)}

\paragraph{Monitoring contract.}
Fix two disjoint finite samples representing \textbf{calibration} and \textbf{deployment}:
\[
  \mathcal{D}_{\mathrm{ref}} = \{(X_i,T_i,Y_i)\}_{i\in\mathcal{I}_{\mathrm{ref}}},
  \qquad
  \mathcal{D}_{\mathrm{live}} = \{(X_j,T_j,Y_j)\}_{j\in\mathcal{I}_{\mathrm{live}}}.
\]
Split $\mathcal{D}_{\mathrm{ref}}$ into training and evaluation subsets
$\mathcal{D}_{\mathrm{ref}}^{\mathrm{tr}}\cup\mathcal{D}_{\mathrm{ref}}^{\mathrm{ev}}$ (held out for in-batch AUUC).
The \textbf{uplift ranker} $\hat\tau:\mathbb{R}^p\to\mathbb{R}$ is fit only on $\mathcal{D}_{\mathrm{ref}}^{\mathrm{tr}}$ and then \textbf{frozen} when scoring $\mathcal{D}_{\mathrm{live}}$.
No refit on LIVE is permitted in the monitoring layer except inside LOCO ablations (always retrain on REF train only).

\paragraph{Batch vs treatment.}
Introduce batch indicator $W\in\{0,1\}$ with $W=0$ on $\mathcal{D}_{\mathrm{ref}}$ and $W=1$ on $\mathcal{D}_{\mathrm{live}}$.
This is \textbf{not} the randomised treatment $T\in\{0,1\}$ used in the uplift experiment.
All FSDS-style batch tests use $(X,W)$ or stacked $(X,Y,W)$; uplift curves use $(Y,T,\hat\tau(X))$ within each batch.

\paragraph{Isomorphism to predictive batch FSDS.}
\begin{center}
\small
\begin{tabular}{@{}lll@{}}
\toprule
\textbf{Predictive FSDS} & \textbf{Two-batch uplift monitor} & \textbf{Object} \\
\midrule
Global loss / calibration on LIVE & $\mathrm{AUUC}_{\mathrm{live}}$, $\Delta_{\mathrm{AUUC}}$ & ranking quality \\
Global MMD$^2$ on $X$ & $\widehat{\mathrm{MMD}}^2(X_{\mathrm{ref}},X_{\mathrm{live}})$ & covariate shift \\
Domain classifier AUC$(X)$ & $\widehat{\mathrm{AUC}}_{\mathrm{dom}}^{\mathrm{RF}}(X;W)$ & same \\
Effective sample size (overlap) & $\mathrm{ESS}_{\mathrm{ovlp}}(\hat e(W\mid X))$ & comparable support \\
LOGO-MMD on feature groups & LOGO-MMD$(g)$ on $X$ & localised $X$ shift \\
Leave-one-group-out on outcome model & LOCO--AUUC$(g)$ & localised ranking drivers \\
Permute--refit on batch (PO-risk) & DRPerm on $(X,Y,W)$ & $Y\mid X$ shift \\
Post-hoc population subset & Domain-quintile AUUC on LIVE & user-slice ranking \\
\bottomrule
\end{tabular}
\end{center}

\subsection{Validity gates (batch geometry)}

\paragraph{Covariate shift on $X$.}
Let $\hat{\phi}$ be an RKHS mean embedding distance or unbiased MMD$^2$ estimator:
\[
  \widehat{\mathrm{MMD}}^2
  = \widehat{\mathrm{MMD}}^2\bigl(\{X_i\}_{i\in\mathcal{I}_{\mathrm{ref}}^{\mathrm{ev}}},
  \{X_j\}_{j\in\mathcal{I}_{\mathrm{live}}}\bigr).
\]
Train a domain classifier $\hat d_W:\mathbb{R}^p\to[0,1]$ on stacked labelled pairs $(X,W)$ and report
$\widehat{\mathrm{AUC}}_{\mathrm{dom}}^{\mathrm{RF}}=\widehat{\mathrm{AUC}}(\hat d_W(X),W)$.

\paragraph{Batch propensity overlap.}
On the stacked sample, estimate $\hat e(w\mid x)=\mathbb P(W=1\mid X=x)$ and define overlap weights
$v_i = 1/(\hat e_i(1-\hat e_i))$ with clipping.
The effective sample size fraction is
\[
  \mathrm{ESS}_{\mathrm{ovlp}}
  = \frac{\bigl(\sum_i v_i\bigr)^2}{\sum_i v_i^2}\cdot\frac{1}{n_{\mathrm{ref}}+n_{\mathrm{live}}}.
\]
Monitoring attributions (LOCO--AUUC, slice tables) require $\mathrm{ESS}_{\mathrm{ovlp}}$ above a policy threshold; otherwise emit \textsc{RefreshRef} before subset rules.

\paragraph{Randomisation check on $T$.}
Within each batch $w\in\{0,1\}$, test treatment balance on $T$ (e.g.\ Pearson $\chi^2$ against design ratio $\pi$).
Failure $\Rightarrow$ invalid uplift curves; halt AUUC interpretation.

\subsection{Layer 1: ranking monitor (AUUC)}

\paragraph{Uplift ranker.}
Let $\mathcal{A}$ be a meta-learner factory (T-/X-/R-learner or registry nuisance models).
Fit $\hat\tau$ on $\mathcal{D}_{\mathrm{ref}}^{\mathrm{tr}}$:
\[
  \hat\tau \leftarrow \mathcal{A}\bigl(\mathcal{D}_{\mathrm{ref}}^{\mathrm{tr}}\bigr).
\]
Scores on evaluation windows:
\[
  \hat\tau_i^{\mathrm{ref}} = \hat\tau(X_i),\ i\in\mathcal{I}_{\mathrm{ref}}^{\mathrm{ev}};
  \qquad
  \hat\tau_j^{\mathrm{live}} = \hat\tau(X_j),\ j\in\mathcal{I}_{\mathrm{live}}.
\]

\paragraph{Area under the uplift curve (AUUC).}
Sort users by $\hat\tau$ descending; let $n_t(k)$, $n_c(k)$ be treated/control counts among the top-$k$ fraction.
With cumulative treated/control response sums $S_t(k)$, $S_c(k)$, define the uplift curve
$U(k)=\frac{S_t(k)}{n_t(k)}-\frac{S_c(k)}{n_c(k)}$ (with standard sklift conventions on ties and empty strata).
Then
\[
  \mathrm{AUUC}_{\mathrm{ref}}
  = \int_0^1 U_{\mathrm{ref}}(k)\,dk
  \quad\text{on}\quad \mathcal{D}_{\mathrm{ref}}^{\mathrm{ev}},
\]
and analogously $\mathrm{AUUC}_{\mathrm{live}}$ on $\mathcal{D}_{\mathrm{live}}$.
The \textbf{two-batch ranking alarm} is
\[
  \Delta_{\mathrm{AUUC}} = \mathrm{AUUC}_{\mathrm{ref}} - \mathrm{AUUC}_{\mathrm{live}}.
\]
Large $|\Delta_{\mathrm{AUUC}}|$ under stable gates indicates deployment ranking degradation relative to REF calibration.

\paragraph{LOCO--AUUC (feature-localised ranking).}
Let $\mathcal{G}=\{g_1,\ldots,g_G\}$ be optional feature groups; if $\mathcal{G}=\emptyset$, each single coordinate defines a group.
For each $g\in\mathcal{G}$, let $X^{(-g)}$ denote features with group $g$ removed.
Retrain $\hat\tau^{(-g)} \leftarrow \mathcal{A}(\mathcal{D}_{\mathrm{ref}}^{\mathrm{tr}}, X^{(-g)})$ and recompute AUUC on both windows:
\[
  \mathrm{drop}_{\mathrm{ref}}^{(g)}
  = \mathrm{AUUC}_{\mathrm{ref}}^{\mathrm{full}} - \mathrm{AUUC}_{\mathrm{ref}}^{(-g)},\quad
  \mathrm{drop}_{\mathrm{live}}^{(g)}
  = \mathrm{AUUC}_{\mathrm{live}}^{\mathrm{full}} - \mathrm{AUUC}_{\mathrm{live}}^{(-g)}.
\]
Define $\mathrm{drop\_gap}^{(g)}=\mathrm{drop}_{\mathrm{live}}^{(g)}-\mathrm{drop}_{\mathrm{ref}}^{(g)}$.
Large $|\mathrm{drop}_{\mathrm{live}}^{(g)}|$ localises LIVE ranking sensitivity to group $g$;
the vector $(\mathrm{drop}_{\mathrm{ref}}^{(g)})_g$ vs $(\mathrm{drop}_{\mathrm{live}}^{(g)})_g$ compared by Spearman $\rho$ distinguishes covariate-driven vs concept-driven reallocation of feature importance.

\paragraph{User slices on LIVE (domain quintiles).}
Train batch propensity on stacked covariates:
\[
  \hat s(x) = \mathbb P(W=1\mid X=x)
  \quad\text{fit on}\quad
  \{(X_i,W_i)\}_{i\in\mathcal{I}_{\mathrm{ref}}^{\mathrm{pool}}\cup\mathcal{I}_{\mathrm{live}}},
\]
where $\mathcal{I}_{\mathrm{ref}}^{\mathrm{pool}}$ is a REF pool (eval $+$ optional train cap).
On LIVE only, form quantile bins $Q_1,\ldots,Q_K$ by $\hat s(X_j)$ and compute
\[
  \mathrm{AUUC}_k = \mathrm{AUUC}\bigl(\mathcal{D}_{\mathrm{live}}\cap Q_k\bigr),\qquad k=1,\ldots,K.
\]
For slices $A,B\subseteq\mathcal{I}_{\mathrm{live}}$, pairwise ranking contrast
$\Delta_{\mathrm{AUUC}}^{A,B}=\mathrm{AUUC}_A-\mathrm{AUUC}_B$; policy ranks pairs by $|\Delta_{\mathrm{AUUC}}^{A,B}|$ for cap/continue rules.

\paragraph{Overlap-support slice.}
Let $\mathcal{L}_{\mathrm{ovlp}}=\{j\in\mathcal{I}_{\mathrm{live}}:\ \hat e(W=1\mid X_j)\in[\ell,u]\}$.
Report $\mathrm{AUUC}_{\mathrm{ovlp}}$ on $\mathcal{L}_{\mathrm{ovlp}}$.
If $\mathrm{AUUC}_{\mathrm{ovlp}}\gg \mathrm{AUUC}_{\mathrm{live}}$, global $\Delta_{\mathrm{AUUC}}$ is partially explained by batch mix off the shared support.

\subsection{Layer 2: uplift level on LIVE slices (effect layer)}

\paragraph{Observed slice uplift.}
On each quintile $Q_k\subseteq\mathcal{D}_{\mathrm{live}}$ with sufficient $T$-support,
\[
  \widehat{\tau}^{\mathrm{obs}}_k
  = \bar Y_{T=1,k}-\bar Y_{T=0,k}.
\]

\paragraph{Model slice uplift level.}
Using the same frozen REF ranker,
\[
  \bar{\hat\tau}_k = \frac{1}{|Q_k|}\sum_{j\in Q_k} \hat\tau(X_j).
\]

\paragraph{Pairwise effect contrasts (LIVE only).}
For quintiles $Q_a,Q_b$,
\[
  \Delta\widehat{\tau}^{\mathrm{obs}}_{a,b}
  = \widehat{\tau}^{\mathrm{obs}}_a - \widehat{\tau}^{\mathrm{obs}}_b,
  \qquad
  \Delta\bar{\hat\tau}_{a,b} = \bar{\hat\tau}_a - \bar{\hat\tau}_b.
\]
Layer~2 policy uses $|\Delta\widehat{\tau}^{\mathrm{obs}}_{a,b}|$ to prioritise budget reallocation \emph{conditional on concept rejection} (below).
Alignment of $\bar{\hat\tau}_k$ with $\widehat{\tau}^{\mathrm{obs}}_k$ across $k$ informs \textsc{Relearn} vs \textsc{Reallocate}.

\paragraph{Layer separation (do not conflate).}
Layer~1 measures \emph{sorting quality} of $\hat\tau$ on LIVE; Layer~2 measures \emph{level} of treatment contrast on LIVE slices.
A slice may exhibit high $\widehat{\tau}^{\mathrm{obs}}_k$ and low $\mathrm{AUUC}_k$ simultaneously; caps follow Layer~1, budget shifts on Layer~2 require stronger concept evidence.

\subsection{Concept axis: PO-risk on stacked batches}

\paragraph{Pseudo-outcome risk.}
Stack $\mathcal{D}=\mathcal{D}_{\mathrm{ref}}\cup\mathcal{D}_{\mathrm{live}}$ with batch labels $W$.
Cross-fit nuisances $\hat\mu(X)$, $\hat e(W\mid X)$; form pseudo-outcomes $\hat\eta=(Y-\hat\mu)(W-\hat e)$ and fit $\hat\tau_{\mathrm{po}}(X)$ by regressing $\hat\eta$ on $X$.
The observed PO-risk statistic is $R=\mathbb E[\hat\tau_{\mathrm{po}}(X)^2]$ (sample mean on $\mathcal{D}$).

\paragraph{Permute--refit test (DRPerm).}
For $b=1,\ldots,B$, permute batch labels $W^{(b)}$, refit nuisances and $\hat\tau_{\mathrm{po}}^{(b)}$, store $R^{(b)}$.
Two-sided Monte Carlo $p$-value:
\[
  p = \frac{1+\sum_{b=1}^B \mathbf 1\{ R^{(b)}\ge R\}}{B+1}.
\]
Reject at level $\alpha$ $\Rightarrow$ interpret large $\Delta_{\mathrm{AUUC}}$ with flat $\widehat{\mathrm{MMD}}^2$ as $Y\mid X$ (and treatment interaction) shift; enable Layer~2 REALLOCATE rules on empirical ATE pairs.

\subsection{Regime classification (two axes)}

\begin{center}
\small
\begin{tabular}{@{}p{0.22\linewidth}cc@{}}
\toprule
 & $Y\mid X$ stable ($p>\alpha$) & $Y\mid X$ shifted ($p\le\alpha$) \\
\midrule
$X$ stable (MMD, AUC flat) &
Layer~1: LOCO + quintile AUUC &
Layer~1 + Layer~2 ATE; PO-LOCO optional \\
\midrule
$X$ shifted (MMD or AUC $\uparrow$) &
LOGO-MMD$(g)$, quintile AUUC, overlap &
Stratify REF; gates then LOCO + L2 \\
\bottomrule
\end{tabular}
\end{center}

If $\mathrm{ESS}_{\mathrm{ovlp}}$ is low, suspend LOCO and slice effect rules; emit refresh/match REF.

\subsection{Monitoring map (formal pipeline)}

\begin{enumerate}
  \item \textbf{Gates:} SRM$(T)$ per batch; $\widehat{\mathrm{MMD}}^2$, $\widehat{\mathrm{AUC}}_{\mathrm{dom}}$; $\mathrm{ESS}_{\mathrm{ovlp}}$.
  \item \textbf{Layer~1 global:} $\mathrm{AUUC}_{\mathrm{ref}}$, $\mathrm{AUUC}_{\mathrm{live}}$, $\Delta_{\mathrm{AUUC}}$.
  \item \textbf{Layer~1 local:} LOCO--AUUC$(g)$; $\mathrm{AUUC}_k$; pairwise $|\Delta_{\mathrm{AUUC}}^{A,B}|$; optional $\mathrm{AUUC}_{\mathrm{ovlp}}$.
  \item \textbf{Concept:} DRPerm $p$ on stacked $(X,Y,W)$.
  \item \textbf{Layer~2 local:} $\widehat{\tau}^{\mathrm{obs}}_k$, $\bar{\hat\tau}_k$; pairwise $|\Delta\widehat{\tau}^{\mathrm{obs}}_{a,b}|$.
  \item \textbf{Policy:} map statistics to $\{\textsc{Reallocate},\textsc{HoldFeature},\textsc{RefreshRef},\textsc{Relearn},\textsc{Monitor}\}$.
\end{enumerate}

\paragraph{Policy mapping (examples).}
Let $\tau_{\mathrm{AUUC}}$ be AUUC SLA, $\delta_{\mathrm{mmd}}$ MMD threshold, $\epsilon_{\mathrm{ess}}$ ESS floor.
\begin{align*}
  \mathrm{ESS}_{\mathrm{ovlp}}<\epsilon_{\mathrm{ess}}
  &\Rightarrow \textsc{RefreshRef}.\\
  \mathrm{AUUC}_{\mathrm{live}}<\tau_{\mathrm{AUUC}}
  &\Rightarrow \textsc{Reallocate}\ \text{(global pause/cap)}.\\
  k^\star=\arg\min_k \mathrm{AUUC}_k
  &\Rightarrow \textsc{Reallocate}\ \text{(cap quintile $Q_{k^\star}$)}.\\
  \arg\max_{a,b}|\Delta\widehat{\tau}^{\mathrm{obs}}_{a,b}|,\ p\le\alpha
  &\Rightarrow \textsc{Reallocate}\ \text{(budget toward higher $\widehat{\tau}^{\mathrm{obs}}$ slice)}.\\
  \widehat{\mathrm{MMD}}^2\ge\delta_{\mathrm{mmd}},\ \mathrm{ESS}_{\mathrm{ovlp}}\ \text{OK}
  &\Rightarrow \textsc{HoldFeature}\ \text{(LOGO-MMD / LOCO top group)}.\\
  p\le\alpha,\ \text{systematic } \bar{\hat\tau}_k\not\approx\widehat{\tau}^{\mathrm{obs}}_k
  &\Rightarrow \textsc{Relearn}\ \text{after REALLOCATE trial}.
\end{align*}

\subsection{Two-batch information flow}

\begin{center}
\begin{tabular}{@{}c|c|c@{}}
\toprule
 & REF & LIVE \\
\midrule
Fit $\hat\tau$ & $\mathcal{D}_{\mathrm{ref}}^{\mathrm{tr}}$ & --- \\
AUUC eval & $\mathcal{D}_{\mathrm{ref}}^{\mathrm{ev}}$ & $\mathcal{D}_{\mathrm{live}}$ \\
LOCO retrain & $\mathcal{D}_{\mathrm{ref}}^{\mathrm{tr}}$ (per drop) & score only \\
Train $\hat s,\hat e$ for slices & $\mathcal{I}_{\mathrm{ref}}^{\mathrm{pool}}\cup\mathcal{I}_{\mathrm{live}}$ & slice LIVE rows \\
Empirical ATE & --- & $Q_k\subseteq\mathcal{D}_{\mathrm{live}}$ \\
PO-risk / DRPerm & \multicolumn{2}{c}{stacked $\mathcal{D}_{\mathrm{ref}}\cup\mathcal{D}_{\mathrm{live}}$} \\
\bottomrule
\end{tabular}
\end{center}

\paragraph{Implementation.}
\texttt{loco\_auuc.loco\_auuc\_monitor} (Layer~1 LOCO);
\texttt{run\_uplift\_subset\_localization} (gates, quintiles, pairwise, rules);
\texttt{run\_uplift\_subset\_benchmark.py} (four REF/LIVE scenarios).
Narrative playbook: \verb|docs/FSDS_UPLIFT_TWO_BATCH_LANDING.md|.


\section{Empirical two-batch scenarios}
% Auto-generated from loco_auuc_benchmark.json — 2026-10-01
\paragraph{Empirical monitoring benchmarks.}
Four REF/LIVE scenarios: Hillstrom (real email uplift); Hillstrom with injected covariate shift on $X$;
synthetic covariate shift ($f_0$ offset on LIVE); synthetic concept shift (new driver $f_3$ on LIVE).
Layer~1 reports Online PFI on $X$ (MMD$^2$, domain RF AUC, overlap ESS); Layer~2 is group LOCO--AUUC
(2 blocks, 6 retrains per learner). AUUC matches \texttt{sklift}; $|\Delta\mathrm{sklift}|=0$ on all runs below.
Monitor overlap ESS on stacked REF/LIVE is high ($\gtrsim 0.7$) in most rows---LOCO attributions are gated in.
\subparagraph{hillstrom.}
FSDS on $X$: MMD$^2$=0.0014, $\widehat{\mathrm{AUC}}_{\mathrm{dom}}^{\mathrm{RF}}=$0.9381, ESS$_{\mathrm{ovlp}}=$0.690.
\begin{center}
\small
\begin{tabular}{lrrrrrr}
\toprule Learner & AUUC$_{\mathrm{ref}}$ & AUUC$_{\mathrm{live}}$ & gap & ESS & $\rho$ & LOCO (s) \\
\midrule
registry\_logistic\_x & -0.0004 & 0.0486 & -0.0490 & 0.982 & 1.00 & 0.50 \\
sklearn\_xlearner & -0.0449 & 0.0154 & -0.0603 & 0.982 & 1.00 & 2.77 \\
registry\_rf\_x & -0.0451 & 0.0017 & -0.0468 & 0.982 & 1.00 & 1.82 \\
tarnet & -0.0227 & -0.0131 & -0.0096 & 0.982 & 1.00 & 7.58 \\
dragonnet & -0.0105 & 0.0463 & -0.0568 & 0.982 & 1.00 & 7.76 \\
\bottomrule
\end{tabular}
\end{center}
\subparagraph{hillstrom\_covariate\_drift.}
FSDS on $X$: MMD$^2$=0.0039, $\widehat{\mathrm{AUC}}_{\mathrm{dom}}^{\mathrm{RF}}=$1.0000, ESS$_{\mathrm{ovlp}}=$1.000.
\begin{center}
\small
\begin{tabular}{lrrrrrr}
\toprule Learner & AUUC$_{\mathrm{ref}}$ & AUUC$_{\mathrm{live}}$ & gap & ESS & $\rho$ & LOCO (s) \\
\midrule
registry\_logistic\_x & 0.0152 & -0.0236 & 0.0388 & 1.000 & 1.00 & 0.43 \\
sklearn\_xlearner & 0.0144 & -0.0181 & 0.0325 & 1.000 & 1.00 & 3.02 \\
registry\_rf\_x & -0.0341 & -0.0093 & -0.0248 & 1.000 & 1.00 & 1.75 \\
tarnet & -0.0241 & 0.0284 & -0.0525 & 1.000 & 1.00 & 5.91 \\
dragonnet & 0.0147 & -0.0164 & 0.0311 & 1.000 & 1.00 & 7.02 \\
\bottomrule
\end{tabular}
\end{center}
\subparagraph{synthetic\_covariate.}
FSDS on $X$: MMD$^2$=0.0029, $\widehat{\mathrm{AUC}}_{\mathrm{dom}}^{\mathrm{RF}}=$0.9999, ESS$_{\mathrm{ovlp}}=$0.753.
\begin{center}
\small
\begin{tabular}{lrrrrrr}
\toprule Learner & AUUC$_{\mathrm{ref}}$ & AUUC$_{\mathrm{live}}$ & gap & ESS & $\rho$ & LOCO (s) \\
\midrule
registry\_logistic\_x & -0.0515 & 0.0231 & -0.0746 & 0.715 & 1.00 & 0.64 \\
sklearn\_xlearner & -0.0160 & 0.0112 & -0.0272 & 0.715 & -1.00 & 6.81 \\
registry\_rf\_x & -0.0277 & 0.0386 & -0.0663 & 0.715 & 1.00 & 2.60 \\
tarnet & -0.0392 & 0.0106 & -0.0498 & 0.715 & 1.00 & 4.56 \\
dragonnet & -0.0566 & 0.0297 & -0.0863 & 0.715 & 1.00 & 5.27 \\
\bottomrule
\end{tabular}
\end{center}
\subparagraph{synthetic\_concept.}
FSDS on $X$: MMD$^2$=0.0015, $\widehat{\mathrm{AUC}}_{\mathrm{dom}}^{\mathrm{RF}}=$1.0000, ESS$_{\mathrm{ovlp}}=$0.893.
\begin{center}
\small
\begin{tabular}{lrrrrrr}
\toprule Learner & AUUC$_{\mathrm{ref}}$ & AUUC$_{\mathrm{live}}$ & gap & ESS & $\rho$ & LOCO (s) \\
\midrule
registry\_logistic\_x & 0.0620 & -0.0328 & 0.0949 & 0.909 & 1.00 & 0.59 \\
sklearn\_xlearner & 0.0088 & -0.0026 & 0.0114 & 0.909 & -1.00 & 6.77 \\
registry\_rf\_x & -0.0820 & -0.0080 & -0.0741 & 0.909 & -1.00 & 2.58 \\
tarnet & -0.0590 & 0.0257 & -0.0847 & 0.909 & -1.00 & 4.50 \\
dragonnet & 0.0261 & -0.0554 & 0.0815 & 0.909 & -1.00 & 5.42 \\
\bottomrule
\end{tabular}
\end{center}
\noindent\emph{Takeaway from the tables.}\quad
AUUC$_{\mathrm{live}}$ moves under both injected covariate shift and concept shift (four scenarios above).
Covariate cases: domain RF AUC on $X$ rises toward 1; AUUC on LIVE often drops or changes sign.
Concept case: AUUC gap is large; LOCO rank correlation $\rho$ flips sign across blocks.
Registry meta-learners give similar AUUC to Dragon-Net on Hillstrom with much faster LOCO runs.


\section{Operational subset insights (summary tables)}
% Comprehensive subset-insight logic for uplift monitoring (tabular FSDS use case).
% No hierarchical feature tree required. \input from uplift_fsds_loco_auuc.tex or standalone.
% Requires: amsmath, booktabs

\section{Subset insights for uplift monitoring}

\paragraph{Goal.}
Global $\mathrm{AUUC}_{\mathrm{live}}$ or $\Delta_{\mathrm{AUUC}}$ flags a problem; \textbf{subset localization} answers \emph{where} on LIVE (which features or which users/slices).
This mirrors distribution-shift FSDS for predictive models, with \textbf{MSE replaced by AUUC} on the ranking layer.
We do \textbf{not} use a hierarchical feature tree here---flat groups, per-feature LOCO, and population quintiles are enough for the uplift tabular use case.

\paragraph{Windows and notation.}
REF (calibration) vs LIVE (deployment); batch $W\in\{0,1\}$; uplift treatment $T$; meta-learner $\hat\tau$ fit on REF train only.
Optional feature \textbf{groups} $\mathcal{G}=\{g_1,\ldots,g_G\}$; if $\mathcal{G}=\emptyset$, LOCO runs \textbf{per feature} (one column dropped at a time).

\subsection{Workflow (order matters when mixture on $X$ is stable)}

\begin{enumerate}
  \item \textbf{Validity.} SRM on $T$; FSDS globals on $X$: $\widehat{\mathrm{MMD}}^2$, $\widehat{\mathrm{AUC}}_{\mathrm{dom}}^{\mathrm{RF}}(X)$; ESS$_{\mathrm{ovlp}}(\hat e(W|X))$ before trusting attributions.
  \item \textbf{Global ranking.} $\mathrm{AUUC}_{\mathrm{ref}}$, $\mathrm{AUUC}_{\mathrm{live}}$, $\Delta_{\mathrm{AUUC}}=\mathrm{AUUC}_{\mathrm{ref}}-\mathrm{AUUC}_{\mathrm{live}}$.
  \item \textbf{LOCO--AUUC} (always when $\Delta_{\mathrm{AUUC}}$ or $\mathrm{AUUC}_{\mathrm{live}}$ fails SLA):
  for each dropped set $g$ (group or single feature $j$), retrain $\hat\tau$ on REF, recompute AUUC on LIVE:
  \[
    \mathrm{drop}_{\mathrm{live}}^{(g)} = \mathrm{AUUC}_{\mathrm{live}}^{\mathrm{full}} - \mathrm{AUUC}_{\mathrm{live}}^{(-g)}.
  \]
  Large $|\mathrm{drop}_{\mathrm{live}}^{(g)}|$ $\Rightarrow$ ranking on LIVE \emph{depends} on that feature subset.
  \item \textbf{Subset localization on LIVE} (after Step 3 when global alert fires):
  \begin{itemize}
    \item \emph{Population slices:} domain-quintile AUUC---partition LIVE by $\hat s(x)=\mathbb P(W{=}1\mid X=x)$; AUUC on each quintile $Q_1,\ldots,Q_5$.
    \item \emph{Pairwise AUUC:} for slices $A,B$, report $\Delta_{A,B}=\mathrm{AUUC}_A-\mathrm{AUUC}_B$; rank pairs by $|\Delta_{A,B}|$ for business rules (cap the worse slice).
    \item \emph{Support slice (optional):} AUUC on LIVE with $\hat e(W|X)\in[\ell,u]$ (overlap band) vs global LIVE.
  \end{itemize}
  \item \textbf{Optional $Y\mid X$ (concept).} PO-risk / DRPerm on stacked REF$\cup$LIVE when Step~1 shows $X$ stable but Step~2 gap is large.
  \item \textbf{Business rules.} Plain sentences: REALLOCATE (cap quintile / overlap-only targeting), HOLD\_FEATURE (high LOCO or LOGO-MMD group), REFRESH\_REF (ESS low), RELEARN (concept after REALLOCATE trial).
\end{enumerate}

\paragraph{When covariate mixture on $X$ is stable ($\widehat{\mathrm{MMD}}^2$ and domain AUC flat).}
User mixture on inputs looks unchanged; a large $\Delta_{\mathrm{AUUC}}$ is then read as \textbf{ranking or concept drift}, not a new population.
\textbf{Do LOCO--AUUC first} (Step 3) to localize \emph{feature}-level drivers; \textbf{then} quintile + pairwise subset localization (Step 4) to localize \emph{who} on LIVE still ranks poorly under the \emph{same} $\hat\tau$.
Add PO-risk (Step 5) if LOCO points to outcome-side change.
Quintiles remain useful because $\hat\tau$ ranking can differ by slice even when $P(X)$ is stable.

\paragraph{When covariate mixture shifts ($\widehat{\mathrm{MMD}}^2$ or domain AUC $\uparrow$).}
Run LOGO-MMD on groups (if $\mathcal{G}\neq\emptyset$): drop group $g$, recompute MMD on $X\setminus g$; large $\mathrm{MMD}_{\mathrm{full}}-\mathrm{MMD}_{\setminus g}$ $\Rightarrow$ group $g$ carries shift on $X$.
Then LOCO--AUUC and quintile AUUC as above (ranking may break on shifted support).

\subsection{Two complementary subset errors (ranking layer)}

\begin{center}
\begin{tabular}{@{}p{0.22\linewidth}p{0.34\linewidth}p{0.36\linewidth}@{}}
\toprule
\textbf{Insight type} & \textbf{Mechanism} & \textbf{Ops read} \\
\midrule
Feature / LOCO--AUUC &
Per-feature or per-group drop, retrain on REF &
Which columns \emph{not} to drop from targeting; which block to refresh \\
Population / quintile AUUC &
AUUC on LIVE slice $Q_k$ &
Which user segment to cap or continue \\
Pairwise AUUC &
$\arg\max_{A,B}|\mathrm{AUUC}_A-\mathrm{AUUC}_B|$ &
Direct A vs B rule (e.g.\ cap $Q_4$, keep $Q_5$) \\
Overlap-band AUUC &
AUUC on comparable-support LIVE &
Global red but core support green $\Rightarrow$ narrow targeting \\
\bottomrule
\end{tabular}
\end{center}

\subsection{Regime summary (what to run first)}

\begin{center}
\small
\begin{tabular}{@{}lll@{}}
\toprule
Signal & Regime & Primary subset tools \\
\midrule
ESS$_{\mathrm{ovlp}}$ low & Mix / incomparable & Refresh REF; no LOCO-driven cuts \\
MMD or domain AUC $\uparrow$, gap $\uparrow$ & Covariate + ranking & LOGO-MMD (if groups), quintile + pairwise \\
MMD flat, gap $\uparrow$ & Mixture stable, ranking/concept & \textbf{LOCO--AUUC first}, then quintile + pairwise; PO-risk \\
\bottomrule
\end{tabular}
\end{center}

\subsection{Main insights (paper-ready)}

\begin{itemize}
  \item $\mathrm{AUUC}_{\mathrm{live}}$ and $\Delta_{\mathrm{AUUC}}$ are \textbf{sensitive to both covariate shift and concept drift}; FSDS on $X$ and PO-risk separate the regime.
  \item \textbf{LOCO--AUUC} localizes error to \emph{features} (groups optional; default per-column if no groups).
  \item \textbf{Domain-quintile AUUC} and \textbf{pairwise comparisons} localize error to \emph{population slices} on LIVE.
  \item When \textbf{$X$-mixture is stable}, interpret large gap via LOCO--AUUC $\rightarrow$ subset quintiles $\rightarrow$ optional PO-risk---not via hierarchical trees.
  \item Output is \textbf{business rules} (REALLOCATE / HOLD\_FEATURE / REFRESH\_REF / RELEARN), not a second model stack.
\end{itemize}

\noindent Implementation: \texttt{loco\_auuc.loco\_auuc\_monitor} (LOCO); \texttt{run\_uplift\_subset\_localization} (quintiles, pairwise, rules).
Empirical tables: \verb|docs/latex/uplift_fsds_benchmark_results.tex|.


\end{document}
